\documentclass[proof]{beamer}
\usetheme{default}

\usecolortheme{rose}
\usepackage[english]{babel}
\usepackage[utf8]{inputenc}
%\usepackage[latin1]{inputenc}
 \usepackage{amssymb}
 \usepackage{latexsym}
 \usepackage{amsmath}
 \usepackage{amssymb}
 \usepackage{tikz}
 \usepackage{tikz-cd}
\usepackage{array}
\usepackage{rotating}
\usepackage{forest}
\usepackage{color}
\usepackage[all,cmtip]{xy}
\usepackage{stmaryrd}

\usepackage{pgfplots}
\usepackage{array}
    \newcolumntype{C}{>{\centering\arraybackslash}c}



\definecolor{qbblue}{RGB}{43,126,128}
%\definecolor{qbblue}{RGB}{43,126,204}
\definecolor{qborange}{RGB}{242,151,36}
\definecolor{qblight}{RGB}{210,210,210}
\definecolor{qbdark}{RGB}{180,180,180}
\definecolor{qbred}{RGB}{184,13,72}

\setbeamercolor{title}{fg=black}
\setbeamercolor{section in toc}{fg=black}
\setbeamercolor{block title}{bg=qblight,fg=black}
\setbeamercolor{frametitle}{bg=qblight,fg=black}
\setbeamercolor{section in head/foot}{bg=black}
\setbeamertemplate{itemize item}{\color{qbdark}$\blacktriangleright$}
\setbeamertemplate{itemize subitem}{\color{qbdark}$\blacktriangleright$}
%\setbeamercolor*{palette tertiary}{bg=black}
\setbeamercolor{local structure}{fg=black}

\usepackage[style=verbose,backend=biber]{biblatex}
\addbibresource{biblio.bib}

\renewcommand{\_}{\rule{.6em}{.5pt}\hspace{0.023cm}}

\newcommand{\bloc}[2]{\begin{block}{#1}\setlength\abovedisplayskip{0pt}#2\vspace{0.1cm}\end{block}}

\setbeamertemplate{navigation symbols}{} 
\addtobeamertemplate{footline}{\hfill{\tiny \insertframenumber}\hspace{2em} \vspace{1em}}

\newcommand{\red}[1]{\textcolor{qbred}{#1}}
\newcommand{\blue}[1]{\textcolor{qbblue}{#1}}
\newcommand{\orange}[1]{\textcolor{qborange}{#1}}

\newcommand{\D}{\mathbb{D}}
\newcommand{\bP}{\mathbb{P}}
\newcommand{\propTrunc}[1]{\lVert #1 \rVert}
\newcommand{\eqspace}{\vspace{0.32cm}}

\newcommand{\SB}{\text{SB}}
\newcommand{\AZ}{\text{AZ}}
\newcommand{\LI}{\Phi}
\newcommand{\B}{\text{B}}
\newcommand{\Et}{\text{\'Et}}
\newcommand{\Aut}{\text{Aut}}
\newcommand{\PGL}{\text{PGL}}
\newcommand{\U}{\text{U}}
\newcommand{\Spec}{\text{Spec}}
\renewcommand{\lim}{\text{lim}}
\newcommand{\M}{\text{M}}

\newcommand{\Alg}{\text{Alg}}
\newcommand{\T}{\mathbb{T}}
\newcommand{\G}{\mathbb{G}}
\newcommand{\Hom}{\text{Hom}}
\newcommand{\Set}{\text{Set}}
\newcommand{\N}{\mathbb{N}}
\newcommand{\PSh}{\text{PSh}}
\newcommand{\Sh}{\text{Sh}}
\newcommand{\J}{\mathbb{J}}
\newcommand{\Type}{\text{Type}}

\newcommand{\spacebefore}{\vspace{0.3cm}}
\newcommand{\spaceafter}{\vspace{-0.2cm}}
\newcommand{\spaceend}{\vspace{-0.4cm}}

\def\mhyphen{{\hbox{-}}}

%\newcommand{\trunc}[1]{|| #1 ||}

\AtBeginSection[]
{
 \begin{frame}<beamer>
 \frametitle{Outline}
 \tableofcontents[currentsection]
 \end{frame}
}

\begin{document}


%\abovedisplayskip=0.0cm
%\abovedisplayshortskip=-0.3cm
%\belowdisplayskip=0.6cm


\title{\red{Sheaves in HoTT \\
with applications to synthetic algebraic geometry}}
\author{Hugo Moeneclaey,\\
\vspace{0.16cm}
Gothenburg Universtity and Chalmers University of Technology\\
}
\date{\blue{HoTT/UF 2026}\\
Aarhus}



\frame{\titlepage}


\frame{\frametitle{Overview}

We present the machinery used to build constructive models for synthetic algebraic geometry and synthetic Stone duality.

\pause

\bloc{What makes it work?}{
\begin{itemize}
\item We want to study some spaces synthetically.
\item Our spaces of interest are build by gluing some basic spaces\footnote{More precisely, they are sheaves on the category of basic spaces.}.
\item These basic spaces have an algebra/geometry duality.
\end{itemize}
}

\pause

%More precisely, you want to work in an $\infty$-topos of sheaf on the opposite of some category of algebras.

\bloc{What prerequisite to use it?}{
\begin{itemize}
\item You don't need technical $\infty$-categories know-how.
\item You do need familiarity with HoTT and its lex modalities.
\end{itemize}
}

}



 %\begin{frame}<beamer>
 %\frametitle{Outline}
 %\tableofcontents
 %\end{frame}
 
 
 
 \section{Introduction to synthetic geometry}
 
 
 \frame{\frametitle{Synthetic geometry}
 
We want to study some spaces synthetically

(e.g.\ smooth manifolds, sequential spaces, schemes).
 
 \pause
 
 \bloc{Method}{
 \begin{enumerate}
 \item Embed the spaces of interest in a topos.
 \pause
 \item Interpret MLTT + UIP in this topos.
 \pause
\item Find \red{new axioms} validated by this interpretation.
\pause
\item Work in MLTT + UIP + \red{these axioms}.
\end{enumerate}
 }
 
 \pause
 
 \bloc{Remark}{
Historical approaches didn't use MLTT.
 }
 
 }
 
 
  \frame{\frametitle{What do we get?}
 
  \bloc{Slogan}{
 \begin{itemize}
 \item All types come with a topology/smooth structure/etc.
 \item All maps are continuous/smooth/etc.
 \end{itemize}
 }
 
 \pause
 
 The spaces of interest are then types satisfying some proposition.

\pause

\bloc{Remark}{
This means we have to abandon LEM and AC.
 }
 
 }

 
 
 \frame{\frametitle{Examples of synthetic geometries}

 
 \renewcommand{\arraystretch}{1.5}
 \begin{center}
\begin{tabular}{|c|c|}
\hline
 Name & Spaces of interest  \\
 \hline
Synthetic differential geometry\footnote{[Lawvere 67], [Kock 77], [Dubuc 79], etc.} & Smooth manifolds \\
\hline
Synthetic domain theory\footnote{[Hyland 91], [Fiore, Plotkin 96], etc.} & Domains \\
\hline
Synthetic topology\footnote{[Escard\'o 04], [Le\v{s}nik 10], etc.} & Sequential spaces \\
\hline
Synthetic algebraic geometry\footnote{[Bleschmidt 17], [Cherubini, Coquand, Hutzler 24], etc.} & Schemes \\
\hline
 {\renewcommand{\arraystretch}{1}
\begin{tabular}{c} 
	Condensed type theory\footnote{[Barton, Commelin 24]}/ \\ 
	Synthetic Stone duality\footnote{[Cherubini, Coquand, Geerlings, Moeneclaey 24]} \\
\end{tabular}} & 
{\renewcommand{\arraystretch}{1}
\begin{tabular}{c} 
	Stone spaces, locally\\ 
	compact Hausdorff spaces \\
\end{tabular} } \\
%Condensed type theory\footnote{[Barton, Commelin 24]}/ & Stone spaces, \\
%Synthetic Stone duality\footnote{[Cherubini, Coquand, Geerlings, Moeneclaey 24]} & compact Hausdorff spaces...\\
\hline
\end{tabular}
\end{center}
 
 }
 
 
 \frame{\frametitle{Synthetic homotopy theory}
 
We want to study $\infty$-groupoids (aka homotopy types) synthetically.
 
\bloc{Method}{
\begin{enumerate}
\item MLTT can be interpreted in $\infty$-groupoids.
\item This validates HoTT (= MLTT + univalence + HITs).
\item Work in HoTT.
\end{enumerate}
}
 
 }
 
 
 \frame{\frametitle{We can do both!}
 
 \bloc{Method}{
 \begin{enumerate}
 \item Embed the spaces of interest in an $\infty$-topos.
 \item Interpret HoTT in this $\infty$-topos.
\item Find \red{new axioms} validated by this interpretation.
\item Work in HoTT + \red{these axioms}.
\end{enumerate}
 }
 
 The spaces of interest are then \red{sets} satisfying some proposition.
 
 \pause

 \bloc{Why do this?}{
 \begin{itemize}
 %\item Some objects of interest have an $\infty$-groupoid structure.
 \item To use some homotopy theory (e.g.\ shape or cohomology).
 \item For the convenience of univalence and HITs.
 \end{itemize}
 }
 
 }
 
 
 \frame{\frametitle{Examples}
 
 \note{We give examples where HoTT/UF features are central. There are HoTT version for many more synthetic geometry.}

 \renewcommand{\arraystretch}{2}
 \begin{center}
 \begin{tabular}{|c|c|}
\hline
 Name & Spaces of interest  \\
 \hline
  \renewcommand{\arraystretch}{1}
 \begin{tabular}{c}
 	Real-cohesive \\
	homotopy type theory\footnote{[Shulman 18]}\\
\end{tabular} & 
 \renewcommand{\arraystretch}{1}
 \begin{tabular}{c}
 	Locally contractible spaces\\
	and their homotopy types\\
\end{tabular} \\
\hline
Synthetic algebraic geometry\footnote{[Bleschmidt 17], [Cherubini, Coquand, Hutzler 24], etc.} & Schemes \\
\hline
 {\renewcommand{\arraystretch}{1}
\begin{tabular}{c} 
	Condensed type theory\footnote{[Barton, Commelin 24]}/ \\ 
	Synthetic Stone duality\footnote{[Cherubini, Coquand, Geerlings, Moeneclaey 24]} \\
\end{tabular}} & 
{\renewcommand{\arraystretch}{1}
\begin{tabular}{c} 
	Stone spaces, locally\\ 
	compact Hausdorff spaces \\
\end{tabular} } \\
\hline
 {\renewcommand{\arraystretch}{1}
\begin{tabular}{c} 
	Simplicial type theory\footnote{[Riehl, Shulman 17], etc.}/ \\ 
	Triangulated type theory\footnote{[Gratzer, Weinberger, Buchholtz 24]} \\
\end{tabular}} & 
Higher categories\\
\hline
\end{tabular}
\end{center}
 
 }
 
 
   \note{MAYBE UNNECESSARY SLIDE?}
 \frame{\frametitle{How to build new synthetic geometries?}
 
 We want to study some spaces synthetically.

\bloc{How to find a suitable $\infty$-topos?}{
\begin{enumerate}
\pause
\item Identify the spaces of interest as gluing of basic spaces.
\pause
\item Find a topology on the category of basic spaces.
\pause
\item Use the corresponding sheaf $\infty$-topos.
\end{enumerate}
} 

\pause

\bloc{How to find the new axioms?}{
Need axioms that are:
\begin{itemize}
\item Validated by the interpretation of HoTT into this topos.
\item Enough to define and study our spaces of interest.
\end{itemize}
}

This is usually very challenging.

}
 
 
 \frame{\frametitle{Algebra/geometry dualities}
 
  \renewcommand{\arraystretch}{1.5}
 \begin{center}
 \begin{tabular}{|c|c|c|}
 \hline
 Name & Algebra & Geometry\\
 \hline
 Zariski duality & Rings & Affine schemes\\
 Stone duality & Boolean algebras & Stone spaces\\
 // & Distributive lattices & Coherent spaces\\
 Gelfand duality & $C^*$-algebras & Compact Hausdorff spaces\\
 \vdots & \vdots & \vdots\\

\end{tabular}
 \end{center}
 
 }
 
 
\frame{\frametitle{Back to our goal}

Let's focus on synthetic algebraic geometry.

\bloc{Goal}{
We want to work synthetically with schemes.
}

Basic spaces are affine schemes.

\pause

\bloc{Core idea}{
Leverage the rings/affine schemes duality to get nice axioms validated by the relevant topoi.
}

\pause

\begin{itemize}
\item \vspace{0cm} [Kock 80] for external duality used in $1$-topoi. 
\item \vspace{0cm} [Bleschmidt 17] for internal duality in $1$-topoi.
\item \vspace{0cm} [Coquand, Höfer, Sattler 26] for internal duality in $\infty$-topoi.
\end{itemize}

%\bloc{Method}{
%\begin{enumerate}
%\item Identifies axioms that holds in the appropriate presheaf models.
%\item Sheafify internally to these models.
%\end{enumerate}
%}

}
 
 
 \section{Presheaf models}
 
 
 \note{TODO this slide is not good}
 \frame{\frametitle{Goal}
 
This section presents results from [Coquand, Höfer, Sattler 26].

\pause

 \vspace{0.25cm}
 We assume given $\T$ an algebraic theory.
 
 \bloc{Idea}{
$\T$ models the duals to our basic spaces.
 }
 
For synthetic algebraic geometry, $\T$ is the theory of rings.

\pause
 
 \bloc{Intermediary goal}{
 Work synthetically in the $\infty$-topos of presheaves on $(\T\mhyphen\Alg_{fp})^{op}$.
 }
 
Need axioms satisfied by the interpretation of HoTT in this topos.
 
% \bloc{Idea}{
% Representables are dual to the algebras, and therefore model the basic building blocks.
% }
 
 }
 
 
 \frame{\frametitle{The generic object}
 
 We have a presheaf sending any $A:\T\mhyphen\Alg_{fp}$ to its underlying set.
 
 It is a $\T$-algebra in $\PSh((\T\mhyphen\Alg_{fp})^{op})$.
 
 \pause
 
 \bloc{Axiom 1 (Generic object)}{
 There is a set $\G$ with a $\T$-algebra structure.
 }
 
 \pause
 
 For synthetic algebraic geometry, we have a ring $R$.
 
 }
 
 
 \frame{\frametitle{The spectrum construction}
 
 \bloc{Definition}{
 A \red{$\G$-algebra} is a $\T$-algebra $A$ with a morphism from $\G$ to $A$.
 }
 
 \pause
 
 \bloc{Definition}{
 Given $A:\G\mhyphen\Alg_{fp}$, we define
 \spacebefore
 \[\Spec(A) = \Hom_{\G\mhyphen\Alg}(A,\G).\]
 \spaceend
 }
 
 \pause
 
 \bloc{Example from synthetic algebraic geometry}{
 Given $P:R[X]$, then
  \spacebefore
 \[\Spec(R[X]/P) = \Hom_{R\mhyphen\Alg}(R[X]/P,R) = \{x:R\ |\ P(x)=0\}.\]
  \spaceend
 }
 
 }
 
 
 \frame{\frametitle{Duality axiom}
 
 \bloc{Axiom 2 (Duality)}{
 Given $A:\G\mhyphen\Alg_{fp}$, the map
 \spacebefore
 \[A \to R^{\Spec(A)}\]
 is an equivalence.
 }
 
 \pause
 
 \bloc{Corollary}{
 The map
  \spacebefore
 \[\Spec : \G\mhyphen\Alg_{fp}\to \Set\]
 is an an embedding.
 }
 
 \pause
 
 \bloc{Definition}{
Sets in the image of $\Spec$ are called \red{affines}. 
 }
 
 }
 
 
 \frame{\frametitle{Duality axiom}
 
  \bloc{Corollary}{
 Given $A,B:\G\mhyphen\Alg_{fp}$, we have that
  \spacebefore
 \[(\Spec(B)\to \Spec(A)) = \Hom_{\G\mhyphen\Alg}(A,B).\]
   \spaceend
 }
 
 \pause
 
 \bloc{Example from synthetic algebraic geometry}{
 \vspace{-0.2cm}
 \[(R \to R) = \Hom_{R\mhyphen\Alg}(R[X],R[X]) = R[X]\]
    \spaceend
 }
 
 }
 
 
 \frame{\frametitle{Choice for affines}
 
% \bloc{Recall}{
% A set $X$ has choice if any surjection to $X$ merely has a section.
% }
 
 We know that representables in a presheaf category are projective.
 
 \pause
 
 \bloc{Axiom 3 (Choice for affines)}{
 Given $A:\G\mhyphen\Alg_{fp}$, the set $\Spec(A)$ has choice.
 }
 
 %This means that any surjection to $\Spec(A)$ merely has a section.
 
 }
 
 
% \frame{\frametitle{MAYBE Digression on qc modules}
 
% }
 
 
% \frame{\frametitle{$\N$ is constant}
 
% $\N$ is interpreted in $\PSh(\T\mhyphen\Alg_{fp})$ as the constant presheaf.
 
% \bloc{Axiom 4 ($\N$ constant)}{
% For any $A:\G\mhyphen \Alg_{fp}$, the map
 
% \[\N\to \N^{\Spec(A)}\]
% is an equivalence.
% }
 
% So family of algebraic objects over $\Spec(A)$ are well-behaved.
 
% \bloc{Proposition}{
% Given $B:\G\mhyphen\Alg_{fp}$, the map
 
% \[(B\mhyphen\Alg_{fp}) \to (\Spec(B)\to \G\mhyphen\Alg_{fp})\]
%is an equivalence.
% }
 
% }
 
 
 \frame{\frametitle{Presheaf models summarized}
 
 \bloc{Theorem [Coquand, Höfer, Sattler 26]}{
 Given $\T$ an algebraic theory, there exists a model of HoTT where:
 \begin{itemize}
 \item There is a $\T$-algebra $\G$.
 \item For all $A:\G\mhyphen\Alg_{fp}$, the map $A\to \G^{\Spec(A)}$ is an equivalence.
 \item For all $A:\G\mhyphen\Alg_{fp}$, the set $\Spec(A)$ has choice.
 %\item For all $A:\G\mhyphen\Alg_{fp}$, the map $\N\to\N^{\Spec(A)}$ is an equivalence.
 \end{itemize}
 }
 
 }
 
 
 
 \section{Sheaf models}


\frame{\frametitle{Toward sheaves}

We want to work with sheaves instead of presheaves.

\bloc{External approach}{
Redo [Coquand, Höfer, Sattler 26] with sheaves.
}

\pause

Forget about higher category theory, we can do it in HoTT!

\pause

\bloc{Internal approach}{
We define sheaves in (HoTT + presheaf axioms).
}

}


 \frame{\frametitle{Internal topologies}
 
 \bloc{Definition}{
A \red{topology} consists of a class $\J$ of finite sums of affines such that $1\in\J$ and $\J$ is stable under $\Sigma$.
 }
 
 \pause
 
 \bloc{Definition}{
 A \red{$\J$-cover} is a map with fibers in $\J$.
 }
 
 \pause
 
 Given $S$ affine, any $\J$-cover over $S$ is merely of the form
 \spacebefore
 \[U_1+\cdots+U_n \to S\]
 with each $U_i$ affine.
 
 \pause

\bloc{Lemma}{
\begin{itemize}
\item Isomorphisms are $\J$-cover.
\item $\J$-covers are stable under pullbacks and composition.
\end{itemize}
}
 
 }
 
 
 \frame{\frametitle{Internal sheaves}

\bloc{Definiton}{
A type $X$ is a \red{$\J$-sheaf} if for all $S\in\J$, the map
\spacebefore
\[X\to X^{\propTrunc{S}}\]
is an equivalence.
}

\pause

\bloc{Lemma}{
Given a $\J$-cover $T\to S$ and a set $X$ that is a $\J$-sheaf, the map
\spacebefore
\[X^S\to \lim(X^T\rightrightarrows X^{T\times_ST})\]
is an equivalence.
}

\pause

\bloc{Remark}{
There is a version for $n$-types based on iterated joins [Williams 25].
\vspace{-0.4cm}
}

 }
 
 \frame{\frametitle{Internal sheaves}

Note that being a $\J$-sheaf is an accessible lex modality.

\bloc{Corollary [Rijke, Shulman, Spitters 20]}{
\begin{itemize}
\item $\J$-sheaves are stable under $\Sigma$ and identity types.
\item $\J$-sheaves are stable under $\Pi_X$ for any type $X$.
\item We have a $\J$-sheafification.
\item The type of $\J$-sheaves is a itself a $\J$-sheaf. 
\end{itemize}
}
  
 }
 

 
 
 \frame{\frametitle{Lex modalities and submodels of HoTT}
 
 Lex modalities correspond to $\infty$-subtopoi.
 
So $\J$-sheaves should form a model of HoTT.

\pause
 
 \bloc{Theorem [Quirin 16]}{
Given an accessible lex modality $L$, we have an interpretation $\llbracket\_\rrbracket$ of HoTT into $L$-modal types where:
\begin{center}
\begin{tabular}{ccc}
$\llbracket \Type \rrbracket$ & $=$ & $\Type_L$ \\
$\llbracket \Sigma_XY \rrbracket$ & $=$ & $\Sigma_{\llbracket X\rrbracket} \llbracket Y \rrbracket$ \\
$\llbracket \Pi_XY \rrbracket$ & $=$ & $\Pi_{\llbracket X\rrbracket} \llbracket Y \rrbracket$ \\
$\llbracket x=_Xy \rrbracket$ & $=$ & $\llbracket x\rrbracket =_{\llbracket X\rrbracket} \llbracket y\rrbracket$\\
$\llbracket \propTrunc{X} \rrbracket$ & $=$ & $L(\propTrunc{\llbracket X\rrbracket})$\\
\end{tabular}
\end{center}
 }
 
 \pause
 
So we have a $\J$-sheaf model of HoTT. 
 
 }
 
 
 \frame{\frametitle{The sheaf axioms: generic object}
 
 What does the $\J$-sheaf model preserve from the presheaf axioms?
 
\pause
 
 \bloc{Definition}{
 A topology $\J$ is \red{subcanonical} if $\G$ is a $\J$-sheaf.
 }
 
 From now on we assume $\J$ subcanonical.
 
\pause
 
 \bloc{Lemma (Generic object in sheaves)}{
 In the $\J$-sheaf model, there is a $\T$-algebra $\G$. 
 }
 
 }
 
 
 \frame{\frametitle{The sheaf axioms: generic object}
 
 So we have $\G$ in the $\J$-sheaf model. Now we want to have $\J$.
 
 \pause
 
 \bloc{Definition}{
A topology is \red{simply defined} if:
\begin{itemize}
\item $\J$ can be defined from (HoTT + there is $\G$ a $\T$-algebra).
\item $\llbracket X\in\J\rrbracket_\J$ implies $L_\J(X\in\J)$.
\end{itemize}
}

From now on we assume $\J$ simply defined.

\pause

\bloc{Lemma}{
In the $\J$-sheaf model, for all $S\in\J$ we have $\propTrunc{S}$.
}
 
}
 
 
 \frame{\frametitle{The sheaf axioms: duality}

 \bloc{Lemma (Duality in sheaves)}{
In the $\J$-sheaf model, duality holds.
 }
 
 }
 
 
 \frame{\frametitle{The sheaf axioms: local choice}
 
 \bloc{Definition}{
 A type $X$ has \red{$\J$-local choice} if given a surjection $Y\twoheadrightarrow X$, there exists a type $\J$-cover $p:U\to X$ with a dotted arrow
 \spacebefore
 \[\xymatrix{
 & Y\ar@{->>}[d]\\
U \ar@{-->}[ru]\ar@{->}[r]_{p} & X 
}\]
\spaceend
 }
 
 \pause
 
 \bloc{Lemma (Local choice)}{
 In the $\J$-sheaf model, affines have $\J$-local choice.
 }
 
 }


\frame{\frametitle{Sheaf models summarized}

 \bloc{Theorem}{
 Assume given:
 \begin{itemize} 
 \item $\T$ an algebraic theory 
 \item $\J$ a subcanonical and simply defined topology.
\end{itemize}
Then there exists a $\J$-sheaf model of HoTT where:
 \begin{itemize}
 \item There is a $\T$-algebra $\G$ such that for all $S\in\J$ we have $\propTrunc{S}$.
 \item For all $A:\G\mhyphen\Alg_{fp}$, the map $A\to \G^{\Spec(A)}$ is an equivalence.
 \item For all $A:\G\mhyphen\Alg_{fp}$, the set $\Spec(A)$ has $\J$-local choice.
 %\item For all $A:\G\mhyphen\Alg_{fp}$, the map $\N\to\N^{\Spec(A)}$ is an equivalence.
 \end{itemize}
 }

}

 
\section{Applications to synthetic algebraic geometry}


\frame{\frametitle{The presheaf model}

%Ring and algebras over them are assumed commutative.

\bloc{Proposition}{
We have a model of HoTT where:
\begin{itemize}
\item There is a ring $R$.
\item For all $A:R\mhyphen\Alg_{fp}$ the map $A\to R^{\Spec(A)}$ is an equivalence.
 \item For all $A:R\mhyphen\Alg_{fp}$ the set $\Spec(A)$ has choice.
\end{itemize}
}

}

\frame{\frametitle{Zariski topology}

\bloc{Definition}{
The Zariski topology consists of the types merely of the form
\spacebefore
\[x_1\ \text{inv}+\cdots+x_n\ \text{inv} \]
where $(x_1,\cdots,x_n)=1$ in $R$.
}

\bloc{Remark}{
$R$ is local if and only if for all $S\in\text{Zar}$ we have $\propTrunc{S}$.
}

}


\frame{\frametitle{The Zariski sheaf model}

\bloc{Proposition}{
We have a model of HoTT where:
\begin{itemize}
\item There is a local ring $R$.
\item For all $A:R\mhyphen\Alg_{fp}$ the map $A\to R^{\Spec(A)}$ is an equivalence.
 \item For all $A:R\mhyphen\Alg_{fp}$ the set $\Spec(A)$ has Zariski-local choice.
\end{itemize}
}

}


\frame{\frametitle{An interlude: sheafifying in a sheaf model}

\bloc{Remark}{
We considered a topology internal to presheaves.

The same story works for a topology internal to sheaves.
}

\pause

\bloc{Remark}{
In Zariski sheaves, affines are stable under finite sums.
}

\pause

Now we work in Zariski sheaves and replace finite sums of affines by affines in the definition of topology.

}


\frame{\frametitle{\'Etale topology}

\bloc{Definition}{
A type $X$ is formally \'etale if for all $\epsilon:R$ nilpotent, the map
\spacebefore
\[X\to X^{\epsilon=0}\]
is an equivalence.
}

\pause

\bloc{Definition}{
The \'etale topology consists of the non-empty formally \'etale affines.
\vspace{-0.4cm}
}

\pause

\bloc{Proposition}{
A type $X$ is an \'etale sheaf if and only if the map 
\spacebefore
\[X\to X^{\propTrunc{\Spec(R[X]/P)}}\]
is an equivalence for all $P$ monic unramifiable.
}

}


\frame{\frametitle{The \'etale sheaf model}

\bloc{Proposition}{
We have a model of HoTT where:
\begin{itemize}
\item There is a local ring $R$ such that any monic unramifiable polynomial in $R$ has a root.
\item For all $A:R\mhyphen\Alg_{fp}$ the map $A\to R^{\Spec(A)}$ is an equivalence.
 \item For all $A:R\mhyphen\Alg_{fp}$ the set $\Spec(A)$ has \'etale-local choice.
\end{itemize}
}

}


\frame{\frametitle{The fppf topology}

\bloc{Definition}{
A affine $\Spec(A)$ is flat if the functor $A\otimes_R\_$ is exact.
}

\bloc{Definition}{
The fppf topology consists of the non-empty flat affines.
}

\pause

\bloc{Remark}{
Given $P$ monic non-constant, $\Spec(R[X]/P)$ is non-empty and flat.

So in fppf sheaves, monic non-constant polynomials have roots.
}

\bloc{Remark}{
Formally \'etale affines are flat, so fppf sheaves are \'etale sheaves.
}

}


\frame{\frametitle{The fppf sheaf model}

\bloc{Proposition}{
We have a model of HoTT where:
\begin{itemize}
\item There is a local ring $R$ such that any monic non-constant polynomial in $R$ has a root.
\item For all $A:R\mhyphen\Alg_{fp}$ the map $A\to R^{\Spec(A)}$ is an equivalence.
 \item For all $A:R\mhyphen\Alg_{fp}$ the set $\Spec(A)$ has flat-local choice.
\end{itemize}
}

}
 
 
 \frame{\frametitle{The four models compared}
 
% They all have duality for finitely presented $R$-algebras. Moreover:
 
 %\begin{center}
 %\begin{tabular}{ccc}
 % & Ring $R$ & Choice for affines \\
 %Presheaves & Arbitrary & Full choice \\
 %Zariski sheaves & Local & Zariski-local choice \\ 
 %\'Etale sheaves & Local and monic unramifiable have roots & \'Etale local choice\\
 %Fppf sheaves & Local and monic non-constant have roots & Flat local choice \\
 %\end{tabular}
 %\end{center}
 
  \begin{center}
 \begin{tabular}{c|c|c}
 & & \\
 Topology & The ring $R$ is & Affines have \\
 & & \\
  \hline
 & & \\
- & arbitrary & full choice \\
 & & \\
Zariski & local & Zariski-local choice \\ 
&  & \\
\'Etale & \begin{tabular}{c} local + monic unramifiable\\ have roots\end{tabular} & \'etale-local choice\\
&  & \\
Fppf & \begin{tabular}{c} local + monic non-constant\\ have roots\end{tabular} & flat-local choice \\
 & & \\
 \end{tabular}
 \end{center}

 
 }
 
 
 \section{Perspective}
 
 \frame{\frametitle{Synthetic Stone duality}
 
 We take $\mathbb{T}$ the theory of boolean algebra. %, or equivalently the theory of boolean rings\footnote{}.

\pause

\bloc{Remark}{
The story we told with finitely presented algebras works similarly with countably presented algebras.
}

\pause

Since boolean algebras are equivalent to boolean rings, 

we can reuse the Zariski, \'etale and fppf topology as is.

\pause

\bloc{Simplifications}{
\begin{itemize}
\item $2$ is the only local boolean ring. 

So in Zariski sheaves $\G = 2$.
\item Any boolean ring is formally \'etale and flat.

So $\text{fppf\ topology} = \text{\'etale\ topology} =  \{S\ |\ \neg\neg S\}$.
\end{itemize}
}
 
}


\frame{\frametitle{A model for synthetic Stone duality}

\bloc{Theorem}{
We have a model of HoTT where:
\begin{itemize}
\item For all $A:\text{Bool}_{cp}$, if $\neg\neg\Spec(A)$ then $\propTrunc{\Spec(A)}$.
\item For all $A:\text{Bool}_{cp}$, the map $A\to 2^{\Spec(A)}$ is an equivalence.
 \item For all $A:\text{Bool}_{cp}$, the set $\Spec(A)$ has local choice.
\end{itemize}
}

}
 
 \frame{\frametitle{Triangulated type theory\footnote{[Gratzer, Weinberger, Buchholtz 24]}}
 
Let $\T$ be the theory of distributive lattice.

Let $\mathbb{L}$ be the generic such lattice.

\bloc{Definition}{
Consider the topology generated by:
\begin{itemize}
\item If $0=_\mathbb{L}1$ we add $\bot$.
\item For all $i,j:\mathbb{L}$ we add $i\leq j + j\leq i$.
\end{itemize}
}

\pause

We should get a model of HoTT with:
\begin{itemize}
\item $\mathbb{L}$ a totally ordered bounded set.
\item Duality for finitely presented distributive lattice over $\mathbb{L}$.
\item Some form of local choice for affines.
\end{itemize}

%See [Williams 25] for preliminary investigations.
 
 }
 
 
 \frame{\frametitle{Summary}
 
 Given \red{spaces of interests} build by \red{gluing basic spaces} such that:
 \begin{itemize}
 \item These basic spaces enjoy a \red{duality} with $\T$-algebras.
 \item The topology $\J$ encoding the gluing can be defined internally.
 \item $\J$ is subcanonical and simply defined.
 \end{itemize}
 
Then we have a model of HoTT where:
\begin{itemize}
\item We have a $\T$-algebra $\G$.
\item $\J$-covers are surjective.
\item We have a \red{duality} between f.p. $\G$-algebras and affines.
\item Affines have $\J$-local choice.
\end{itemize}

\pause

In SAG and SSD, these axioms are enough for substantial results.
  
 }
 
% \frame{\frametitle{Maybe an overview of SSD and SAG results}}
 

\end{document}